\section{攀登无尽的雪山：AI 创业的心路}

\subsection{从"向延绵而未知的雪山前进"到"Beginning of Infinity"}

杨植麟是月之暗面（Moonshot AI）创始人兼 CEO。时隔一年（2025年7月），再次接受访谈时，他用《无穷的开始》（\textit{The Beginning of Infinity}）的哲学来描述自己的感受：

\begin{importantbox}{问题不可避免，但问题可以解决}
David Deutsch 在《无穷的开始》中提出两句核心命题：
\begin{enumerate}
    \item 问题是不可避免的（Problems are inevitable）
    \item 问题是可以解决的（Problems are soluble）
\end{enumerate}
研发 AI 的过程恰好体现了这一哲学：解决了强化学习的很多问题，就面临评估、验证等新问题；解决了预训练效率问题，又面临 Agent 泛化问题。每次解决一个问题，技术就再攀登几百米，同时看到新的问题。
\end{importantbox}

杨植麟认为，AGI 可能不是某一级台阶，而是一个方向。今天的 AI 在很多领域已经比 99\% 的人类做得更好——在某种意义上可以视为部分 AGI。但很难说某个时间点突然达到 AGI，它更像是持续攀登的过程。

\begin{knowledgebox}{雪山可能没有尽头}
杨植麟表达了一个乐观的观点：他希望这座雪山一直没有尽头——"所谓的 Beginning of Infinity 就是这个意思，它可能是一座无限的山"。而且到某个阶段，可能不再是人类在爬，而是用 AI 来攀登——比如用 K2 模型来参与模型训练、数据处理相关的工作。AI 正在成为"人类文明的放大器"。
\end{knowledgebox}

\subsection{本章小结}

AI 创业是一个持续解决问题、又发现新问题的过程。AGI 不是一个终点，而是一个方向。AI 正在从工具演变为研发过程中的参与者。


